\Needspace{10\baselineskip}
\section{Silima as a Broadcasting Platform}

Silima is designed to operate as a locally hosted gameplay
broadcasting and interaction platform. Although its
implementation differs from commercial streaming services,
its intended functionality is comparable in several
important respects.

\subsection{Technical Functionality}

Conventional gameplay streaming services, such as Twitch
and YouTube Live, allow a host to execute a video game
and transmit its audiovisual output to viewers over
a network.

Silima is designed around a similar model. A video game
executes on a host computer, while connected participants
interact with the game through a locally hosted website.
Silima's planned broadcasting functionality will allow
participants to view the resulting gameplay through
their web browsers.

Unlike conventional streaming services, Silima also
allows connected participants to collectively influence
gameplay by submitting or voting on controller inputs.
These inputs are processed by the host application and
forwarded to the game through RetroArch.

The primary technical distinctions are:
\begin{itemize}
    \item Silima operates over a local area network
          rather than requiring an internet connection.
    \item Gameplay is executed entirely on the host
          computer rather than on participant devices.
    \item Participants may collectively influence
          gameplay through a shared controller interface.
    \item Silima does not require a third-party
          broadcasting or streaming service.
\end{itemize}

We believe that operating over a private network,
rather than the public internet, does not fundamentally
change the technical nature of transmitting gameplay
to an audience.

However, this technical similarity does not necessarily
mean that Silima is legally equivalent to a commercial
streaming platform or that publisher permissions
automatically extend to it.

\subsection{Broadcasting Versus Game Distribution}

An important distinction exists between broadcasting
gameplay and distributing the software used to
execute a game.

When gameplay is streamed through a conventional
broadcasting service, viewers receive audiovisual
content rather than the executable game software.

Silima follows this same general distinction.
Games are executed locally through RetroArch on
the host computer. Participating devices do not
receive or execute the game software itself.

Silima does not provide or distribute:
\begin{itemize}
    \item Commercial game executables or ROM files.
    \item Console firmware or BIOS files.
    \item Copies of commercial games for installation
          on participant devices.
\end{itemize}

Gameplay broadcasting may still involve the
transmission of copyrighted audiovisual material.
The distinction between broadcasting and software
distribution does not eliminate copyright
considerations associated with gameplay footage,
music, or other protected content.

The host remains responsible for obtaining and
using game software in accordance with applicable
laws and licensing requirements.

\subsection{Interpretation of Publisher Policies}

Many video game publishers have established policies
permitting individuals to share gameplay footage
through streaming services, video-sharing platforms,
and other forms of online broadcasting.

These policies recognize gameplay broadcasting as
an established activity, although the permissions
and restrictions vary between publishers.

We believe that Silima's intended broadcasting
functionality is sufficiently similar to conventional
gameplay streaming to warrant consideration under
these existing policies.

In particular, Silima:
\begin{itemize}
    \item Executes games on a host computer while
          transmitting gameplay to an audience.
    \item Allows participants to interact with and
          influence the resulting gameplay.
    \item Is intended for noncommercial classroom
          activities.
    \item Does not distribute executable game
          software to participating devices.
\end{itemize}

Nevertheless, publisher policies may restrict
broadcasting to particular platforms, individuals,
types of content, or methods of distribution.
Some may also require commentary or other
original contributions.

Consequently, we do not assume that permission
to broadcast gameplay through one service
automatically authorizes its transmission
through Silima.

Instead, we interpret the similarities between
Silima and established broadcasting services as
a basis for evaluating individual publisher
policies. The following sections examine those
policies and their applicability to Silima.
